\documentstyle[%


righttag%


,11pt%


,amssymb%


,russian%


,izvru%


]{amsart}





\newcommand{\tl}[3]{\medskip\noindent{\bf#1}\ #2 \dotfill #3 \par} 


\pagestyle{empty}


\begin{document}


%\tableofcontents


\par


\vskip 2cm


\par


\bigskip


\bigskip


\bigskip


\par


\centerline{\Large СОДЕРЖАНИЕ}


\bigskip


\bigskip


%%%%%%%%%% Use \newline %%%%%%%%%%%





\tl{Бояршинова А.В.}


{О пространстве существенных инфинитезимальных деформаций}{3}





\tl{Витюк А.Н.} 


{Существование решений дифференциальных включений 


с частными\newline производными дробных порядков}{13}





\tl{Володин А.И.} 


{Сходимость взвешенных сумм случайных элементов 


в банаховых\newline пространствах типа $p$}{20}





\tl{Егишянц С.А., Островский Е.И.} 


{Асимптотическое поведение случайных полей 


на\newline последовательности расширяющихся множеств}{26}





\tl{Коcов А.А.} 


{О глобальной ycтойчивоcти


неавтономныx cиcтем.{\rm II}}{33}





\tl{Секерин А.Б.} 


{Теорема о носителе для комплексного 


преобразования Радона}{43}





\tl{Сосов Е.Н.} 


{О геодезических отображениях 


специальных метрических пространств}{46}





\tl{Тронин С.Н.} 


{О Морита-контекстах, возникающих при изоморфизме


некоторых подколец \newline колец эндоморфизмов}{50}





\tl{Якупов М.Ш.} 


{Каноническое описание и обращение операторов 


в дискретном\newline пространстве. {\rm I}}{59}








\bigskip


{\centerline {Краткие сообщения}}


\bigskip





\tl{Аксентьев Л.А., Бильченко Г.Г.} 


{К обратной задаче для интегралов 


Кристоффеля--\newline Шварца}{72}





\tl{Ратинер Н.М.} 


{Свойства классов бордизмов оснащенных


подмногообразий фредгольмовой\newline структуры}{77}











\vfill


\noindent\raisebox{2pt}{\copyright} Казанский госудаpственный унивеpситет


\end{document}


